针对手机银行转账场景中日益突出的语义欺诈问题，传统风控系统主要依赖交易金额、设备状态、登录地点和交易频次等结构化字段进行判断，对冒充客服、虚假投资、公检法威胁等社交工程类风险的识别能力相对有限。此类欺诈往往发生在用户本人设备和本人账户的正常操作过程中，交易参数未必表现出明显异常，风险更多隐藏在转账意图、用户补充说明和操作过程细节之中。围绕这一问题，本文设计并实现了一套基于 AI Agent 与模型上下文协议（Model Context Protocol，MCP）的手机银行语义欺诈风控原型系统，尝试将语义理解、微行为感知、规则治理和审计追踪纳入同一条转账风控链路。

系统总体上采用 HIRD 分层架构，将处理流程划分为感知集成层、认知决策层、规则治理层和业务执行层。感知集成层面向移动端转账过程，采集输入时长、异常停顿、内容修改、页面停留和操作节奏等微行为特征，并与交易金额、收款对象、设备信息和交易备注共同封装为风险快照。认知决策层利用 AI Agent 对交易文本和用户补充说明进行语义分析，并结合检索增强生成（Retrieval-Augmented Generation，RAG）机制引入欺诈案例和规则片段，降低单纯依赖模型内部知识进行判断的不稳定性。规则治理层不直接采纳大模型的开放式输出，而是将语义风险分、结构化规则分和场景调节因子融合为可解释的风险结果，并通过阈值路由生成放行、二次核验或拦截等确定性处置指令。业务执行层则负责交易状态流转、人工复核、日志记录和敏感信息脱敏，使 Agent 主要承担分析任务，而不直接控制资金类操作。

在工程实现方面，系统后端基于 FastAPI 完成异步接口编排和冷热路径分流，移动端原型用于展示转账提交、风险提示、二次核验和复核反馈等流程。MCP 被用于封装风险计算、审计记录和业务执行等原子工具，使工具调用具备统一的输入输出格式、幂等校验和 Trace ID 追踪能力。测试阶段围绕 66 个后端自动化用例和 54 个核心场景样本进行验证。结果显示，系统在构造样本集上的精确匹配率为 92.6\%，欺诈样本召回率为 100\%，正常交易直接误拦截率为 0\%。性能测试中，预检热路径平均响应时间为 43.60 ms，P95 响应时间为 78.39 ms；二次核验路径平均响应时间为 2360.31 ms，P95 响应时间为 2767.42 ms。实验结果说明，本文原型系统能够在当前测试条件下完成语义欺诈识别、分级处置和审计追踪闭环，同时也表明其效果仍受样本规模、真实业务数据和外部模型调用稳定性的限制，后续还需要在更复杂的真实场景中继续验证和优化。
